\section{Data, Ethics, and Practical Considerations}

\subsection{Data Readiness}

The primary dataset is the eICU Collaborative Research Database Demo (Pollard et al., 2018), a freely available subset of a larger ICU database hosted on PhysioNet. The full eICU database covers 73,718 patients from 208 US hospitals. The demo subset used in this project contains 2,520 patients across 186 hospitals, with information on demographics, APACHE physiology scores (heart rate, blood pressure, creatinine, GCS components, temperature, and similar clinical measurements), lab results, diagnoses, and whether each patient survived their hospital stay. The multi-hospital structure is what makes this dataset suitable for the project  each hospital or geographic region acts as an independent clinical environment with its own patient mix and clinical practices, which is exactly the kind of cross-context variation the framework needs to work with.
Availability and licensing. The demo subset can be downloaded from PhysioNet without any application or credentialing process (https://physionet.org/content/eicu-crd-demo/2.0/). The full database requires a free online training course (CITI "Data or Specimens Only Research") and a signed data use agreement. The licence is specific to PhysioNet  it allows research use but does not permit redistribution or commercial use without a separate arrangement. This is not a Creative Commons licence, so the usual CC assumptions about sharing and reuse do not apply. Existing benchmark code from the YAIB framework (Yèche et al., 2023) provides ready-made data processing pipelines and baseline models for ICU mortality prediction, which means the project can focus on the new diagnostic framework rather than spending time on data engineering.
Data Readiness Level. On the Lawrence Framework scale, the eICU demo sits at roughly DRL-3  the data is curated and structured but needs task-specific work before modelling. In practice, this means joining several CSV files (patient records, physiology scores, hospital metadata, and outcomes), handling missing values, and recoding certain fields like age. After this preprocessing, the data reaches DRL-5 and is ready for the mortality prediction task. Since the project is a prediction task needing many rows rather than many columns, and the data is already structured rather than raw, the balance of time favours analysis over cleaning.
Quality and limitations. The data quality is high  these are real records from working ICUs using standardised electronic health record systems. The APACHE physiology variables are measured in a consistent way across hospitals, which reduces the risk that differences in model performance across sites are just measurement artefacts. There are three main limitations. First, every hospital is in the United States, so the framework cannot be tested on non-US healthcare systems where coding practices, treatment approaches, and patient demographics may differ significantly. Second, the demo subset only has 10 to 40 patients per hospital (median 12), which is too few for hospital-level modelling  the project handles this by grouping hospitals into US regions (Midwest, South, West, Northeast), giving four contexts with 140 to 807 patients each. The full database would not have this constraint. Third, the patient population is 80% Caucasian, predominantly older, and majority male, which limits how well the framework's thresholds can be tested for fairness across different demographic groups.
4.2 Ethical Red-Teaming
Fairness and representation. The framework's deployment thresholds are calibrated using data that reflects a specific slice of the US patient population  80% Caucasian, with African American patients at 11%, Hispanic at 6%, and Asian at 1.5%. Thresholds set on this data may not work properly for populations that look different. The smallest regional group (Northeast, 140 patients) has the highest mortality rate (13.6%, compared to 7.1% in the Midwest), which shows that patient characteristics do vary meaningfully between contexts. The risk is that a framework built to detect generalisation failure could itself give misleading reassurance to the groups who are most at risk of being poorly served by clinical models.
Privacy and identity. The data is de-identified, but some hospitals in the demo have as few as 10 patients. When you compute statistics about small subgroups  which divergence scoring requires  the combination of unusual demographic details with specific clinical features could potentially allow someone to be identified, especially if the data were linked with other sources. Under UK GDPR, de-identified data can still count as personal data if re-identification is realistically possible. Under the DUAA 2025, if the framework's recommendations were used in clinical practice, patients would have the right to have a human review any automated decision that affects them.
Harm and misuse. The framework gives categorical outputs  deploy, restrict, or do not deploy. There is a risk that people treat these as definitive answers rather than estimates. A model that just barely passes the threshold might be deployed with more confidence than the evidence supports, especially in organisations that are under pressure to adopt AI tools. The opposite problem also exists: a false alarm could pull a useful model out of service unnecessarily. There is also an ecological fallacy issue  the divergence score describes a population, but predictions are made about individual patients. A low divergence score for a region does not mean that every patient in that region is well-represented in the training data.
Data governance. The PhysioNet licence allows research use but requires that the work stays within proper governance boundaries. Individual patient records should not be extracted, published, or shared beyond the data use agreement. If the framework were ever deployed in a UK clinical setting, the DUAA 2025 rules on automated decision-making would apply  meaning a human must be able to review and override any automated recommendation, and the reasoning behind it must be explainable.
4.3 Technical Mitigations
For fairness: Break down all threshold results by demographic subgroups where sample sizes allow. Include a clear statement of which populations were represented in the calibration and which were not. Do not claim that thresholds set on predominantly white US hospital data will apply to other populations.
For privacy: Keep all work within PhysioNet's data governance rules. Do not report statistics for any subgroup with fewer than 10 people. Calculate divergence scores using aggregate distributions, not individual records. Do not link eICU data with any outside datasets.
For automation bias: Show the actual divergence score and its uncertainty range alongside the categorical recommendation, so practitioners can see how close to the boundary they are. Require a documented human decision before any model is deployed, paused, or withdrawn based on the framework. Make it clear in all outputs that the recommendation is an estimate, not a guarantee.
For feature exclusion equity: Before removing shift-sensitive features and retraining, check whether dropping those features hurts performance for specific demographic groups more than others. If it does, use targeted recalibration for those groups instead of removing the features entirely.
For data governance: Follow all PhysioNet data use terms. If the framework were used in practice in the UK, put DUAA 2025 safeguards in place  a human must be able to review any automated recommendation, and the basis for the recommendation must be recorded and auditable.
